\documentclass[11pt,a4paper]{article}
\usepackage[T1]{fontenc}
\usepackage[utf8]{inputenc}
\usepackage{amsmath,amssymb,xcolor}
\usepackage[margin=25mm]{geometry}
\usepackage[colorlinks=true,urlcolor=blue]{hyperref}
\definecolor{nexorablue}{HTML}{2459B5}
\setlength{\parindent}{0pt}
\setlength{\parskip}{7pt}
\setlength{\emergencystretch}{2em}
\renewcommand{\familydefault}{\sfdefault}
\begin{document}

{\small\bfseries\color{nexorablue} NEXORA CPGE}\par
{\small MP / NEX-TOP-05 / Énoncé}\par
{\LARGE\bfseries Un minimum sur une droite non bornée}\par
\vspace{0.25em}\hrule\vspace{0.6em}

Sur la droite $D=\{(x,y):x+2y=3\}$, on étudie $f(x,y)=(x-2)^2+3y^2$. Établir l’existence d’un minimum en se ramenant à un compact, puis calculer sa valeur et le point où il est atteint.


\end{document}
